\makeatletter
\def\input@path{{../styles/}{../../styles/}{../../../styles/}{../}{../../}{../../../}}
\makeatother


\documentclass{ee102_pset}
% macros.tex - Course meta information
\renewcommand{\course}{EE 102} % with a space
\renewcommand{\coursetitle}{Signal Processing and Linear Systems}
\renewcommand{\instructor}{Ayush Pandey}
\renewcommand{\student}{Name: }

\renewcommand{\semester}{Fall 2025}
\date{\semester} % this sets the LaTeX date field safely

% Problem set number
\renewcommand{\psetnum}{1}

% Release / due — use renewcommand because package provides empties
\renewcommand{\releasedate}{August 20, 2025}
\renewcommand{\duedate}{August 27, 2025}

\input{preamble.tex}

% Assignment info
\author{\rule{3cm}{0.4pt}} % Name placeholder
\submitdate{\rule{3cm}{0.4pt}} % Submission date placeholder
\problemset{Problem Set \#7: Applications of Fourier Series}
\renewcommand{\psetnum}{7}
\renewcommand{\releasedate}{October 12, 2026}
\renewcommand{\duedate}{October 21, 2026}
\shorttitle{Problem Set \#7}
\hypersetup{pdftitle={EE 102 Problem Set 7: Applications of Fourier Series}}

\begin{document}
\problem{1}
[Adapted from Oppenheim and Willsky, Problem 3.62]

An ideal full-wave rectifier passes the positive half-cycles of an AC input unchanged and reverses the polarity of the negative half-cycles. So, its output has a DC component and AC ripple. This is usually the first stage of constructing a DC
power supply. We are \textbf{not} modeling the next stages that use a smoothing filter and we are not modeling the diode voltage drops.
Answer the following questions about this system. Show your reasoning and include your code for the computational parts.

\problempart[10 points] First step is to implement the modeling work we have done in this class. Sketch a block diagram with the AC input $x(t)$ and unipolar output $y(t)$, and write $y(t)$ as a function of $x(t)$. For a specific $x(t)=\cos(t)$, sketch the
input and output and find their fundamental periods. 

\problempart[5 points] Show that the rectifier system is nonlinear.

\textbf{Note:} For the parts that follow, use $x(t)=\cos(t)$ and its rectified output $y(t)$. We denote the fundamental period of the output as $T_y$ (you will find it in part (a)). The
fundamental angular frequency of $y(t)$ is $\omega_y=2\pi/T_y$, and its complex Fourier series coefficients are denoted as $a_k$.

\problempart[25 points] Compute the Fourier series coefficients $a_k$ of the rectified output.

Note that before doing the integration, you can observe the symmetry of the signal to predict whether $a_k$ will be real or imaginary and even or odd as a sequence in $k$ (we discussed these properties in class at length!).

Hint: You may use a period from
$[-\pi/2,\pi/2]$, so the integration bounds are
\[
a_k=\frac{1}{T_y}\int_{-\pi/2}^{\pi/2}
y(t)e^{-jk\omega_y t}\,dt.
\]
In this time period, what is $y(t)$ equal to?

\problempart[20 points] Let's compute the same coefficients but this time using the multiplication property of Fourier series. 

Note that if two signals $u(t)$ and $v(t)$ are represented using the same reference
period $T$ and $\omega_0=2\pi/T$, with coefficients $p_k$ and $q_k$, then
their product $z(t)=u(t)v(t)$ has coefficients
\[
c_k=\sum_{m=-\infty}^{\infty}p_m q_{k-m}.
\]
Thus multiplication in time corresponds to convolution of the coefficient
sequences (week 7 lecture 1). This property does not require the rectifier to be linear!

Important hint: Write the rectified output as
\[
y(t)=\cos(t)\operatorname{sgn}(\cos(t)),
\]
where $\operatorname{sgn}(v)$ is $1$ for $v>0$, $-1$ for $v<0$, and $0$
for $v=0$. Use the common reference period $T=2\pi$ and $\omega_0=1$
for both factors in this part.
\begin{itemize}
  \item Find the coefficients $p_k$ of $\cos(t)$ using Euler's formula.
  \item Compute the coefficients $q_k$ of $\operatorname{sgn}(\cos(t))$, including $q_0$, by hand using the analysis integral over $[-\pi,\pi]$.
  \item Apply the multiplication property to obtain the output
  coefficients $c_k$. Simplify the convolution sum using the nonzero $p_k$.
  \item Compare your result with part (c). Note that the reference period in the computation above is twice the fundamental period of the output, so the coefficients you find will be related
  as $a_k=c_{2k}$. 
\end{itemize}

\problempart[10 points] Now, let us try to combine the Fourier series parts to reconstruct the output waveform computationally.
Define the partial sum, just like we did in the lecture notes:
\[
y_N(t)=\sum_{k=-N}^{N}a_k e^{jk\omega_y t}.
\]
Using computer programming, plot $y(t)$ together with $y_N(t)$ for
$N=1,3,10$ over $[-\pi/2,\pi/2]$. Include the DC term and the first $N$
positive and negative harmonics, for $2N+1$ complex-series terms in total.
Describe how the approximation changes as $N$ increases, especially near
the corners of the rectified waveform.

\problempart[20 points] We studied that Fourier Series reconstruction approaches the original waveform, in energy computation, as $N\to\infty$. So, let's put this to test. 
Recall Parseval's theorem:
\[
E_y=\int_{-\pi/2}^{\pi/2}|y(t)|^2\,dt
=T_y\sum_{k=-\infty}^{\infty}|a_k|^2
=T_y\left(|a_0|^2+2\sum_{k=1}^{\infty}|a_k|^2\right).
\]
The final equality follows from the conjugate symmetry of a real signal
coefficients. For this real, even output, the coefficients themselves are
real, so $|a_k|^2=a_k^2$.
\begin{itemize}
  \item Compute $E_y$ exactly by integrating the squared waveform.
  \item Compute an approximation to $E_y$ by numerically summing squared
  coefficient magnitudes for $k=-K,\ldots,K$. Report your cutoff $K$ and explain your findings.
  \item Find the DC level $a_0$. Define the AC ripple as
  $y_{\mathrm{ripple}}(t)=y(t)-a_0$, and use the coefficient formulas
  \[
  E_{\mathrm{DC}}=T_y|a_0|^2,\qquad
  E_{\mathrm{ripple}}=2T_y\sum_{k=1}^{\infty}|a_k|^2
  =E_y-E_{\mathrm{DC}}
  \]
  to compute the DC and ripple energies over one output period. What
  percentage of $E_y$ is in each component?
\end{itemize}

\problempart[10 points] Finally, let us establish that the error energy really decreases as more harmonics are added. So, let's apply Parseval's theorem to the error signal now.

Define $e_N(t)=y(t)-y_N(t)$ and its
one-period error energy
\[
E_N=\int_{-\pi/2}^{\pi/2}|e_N(t)|^2\,dt.
\]
Then, Parseval's theorem gives the explicit formula
\[
E_N=E_y-T_y\sum_{k=-N}^{N}|a_k|^2
=T_y\sum_{|k|>N}|a_k|^2.
\]
Using the exact $E_y$ from part (f) and the finite retained-coefficient
sum, compute $E_N$ for $N=1,\ldots,20$ and plot $E_N$ versus $N$.

\end{document}
